\section*{Chapitre \ref{ch:m2:short-term-interest-rate-futures} --- Contrats à terme de taux à court terme}

\begin{solution}{exo:m2:short-term-interest-rate-futures:1}
$100 - 95.965 = 4.035\%$. Une hausse de 3.5 points de base correspond à 7 pas de cotation d'un demi-point de base~;
sur 200 contrats, le gain vaut $3.5 \times 25 \times 200 = \text{USD}~17\,500$.
\end{solution}

\begin{solution}{exo:m2:short-term-interest-rate-futures:2}
17 mars, 16 juin, 15 septembre et 15 décembre 2027. Le \omterm{def:m2:short-term-interest-rate-futures:imm}{trimestre de référence} du contrat de juin
2027 va du 17 mars (inclus) au 16 juin 2027 (exclu).
\end{solution}

\begin{solution}{exo:m2:short-term-interest-rate-futures:3}
$100 - 96.1025 = 3.8975\%$. Une position longue perd lorsque les taux montent~: $2 \times 41.67
\times 100 = \text{USD}~8\,334$.
\end{solution}

\begin{solution}{exo:m2:short-term-interest-rate-futures:4}
Seuls trois des 31 jours d'octobre portent le nouveau taux, si bien que le multiplicateur vaut $31/3$~:
une erreur d'un point de base sur le prix déplace le taux implicite de 10.3 points de base, soit 41~\% d'une
hausse. Utiliser le contrat de novembre, dans lequel le taux postérieur à la réunion s'applique chaque jour
(il n'y a pas de réunion en novembre)~: son cours donne directement ce taux.
\end{solution}

\begin{solution}{exo:m2:short-term-interest-rate-futures:5}
$[(1 + 0.0395 \times 91/360)(1 + 0.0405 \times 91/360)(1 + 0.0410 \times 91/360)(1 +
0.0412 \times 91/360) - 1]/(4 \times 91/360) = 4.1178\%$.
\end{solution}

\begin{solution}{exo:m2:short-term-interest-rate-futures:6}
Les deux bourses ne désignent pas leurs contrats de la même façon. Le contrat à trois mois sur le SOFR
«~décembre~» capitalise de mi-septembre à mi-décembre~: la
période \emph{qui se termine} en décembre. Le contrat «~décembre~» à trois mois sur le SONIA
s'accumule à partir du troisième mercredi de décembre jusqu'à mi-mars~: la
période \emph{qui commence} en décembre. La couverture est décalée d'un trimestre~: il faut lire la
période de référence dans les règles de chaque contrat, et jamais le mois figurant dans son nom.
\end{solution}

\begin{solution}{exo:m2:short-term-interest-rate-futures:7}
À 1~\%~: $\frac12 \times 0.0001 \times 5 \times 5.25 = 13.1$ points de base à cinq ans~;
$\frac12 \times 0.0001 \times 10 \times 10.25 = 51.25$ à dix ans. À 1.2~\%, les deux montants sont multipliés par
$1.44$~: 18.9 et 73.8 points de base, soit une hausse de 5.8 et 22.55.
\end{solution}

\begin{solution}{exo:m2:short-term-interest-rate-futures:8}
Le contrat d'octobre n'a que trois jours au nouveau taux, si bien que sa probabilité implicite
est dix fois plus exposée au bruit de prix que celle de novembre~: un écart
d'un point de probabilité ne représente qu'une fraction de tick dans le prix
d'octobre. Les deux nombres ne portent pas sur la même chose~: la
moyenne d'octobre contient aussi trois jours du passage de fin de mois, quel qu'il soit. Aucune
de ces deux différences n'est un arbitrage~: la lecture de novembre est la plus fiable.
\end{solution}

\begin{solution}{pb:m2:short-term-interest-rate-futures:1}
\textbf{1.} $100 - 96.1025 = 3.8975\%$.
\textbf{2.} Neuf (du 1 au 9 décembre) au taux précédent, vingt-deux au nouveau.
\textbf{3.} $(31 \times 3.96 - 9 \times 3.8975)/22 = 3.9856\%$.
\textbf{4.} $(3.9856 - 3.8975)/0.25 = 35.2\%$.
\textbf{5.} Un point de probabilité vaut 0.25 points de base du nouveau taux sur 22 des
31 jours~: $0.25 \times 22/31 = 0.18$ points de base du prix de décembre.
\textbf{6.} $(31 \times 3.96 - 0.10 - 9 \times 3.8975)/22 = 3.9810\%$~: 33.4~\%.
\textbf{7.} 3.9765~\%~: 31.6~\%.
\textbf{8.} $0.10/22 \times 100/0.25 = 0.18$ points de pourcentage par point de base de
\omterm{def:m2:money-markets:turn}{passage de fin de période}.
\textbf{9.} Un contrat d'un mois fait la moyenne des jours calendaires de son mois~; le
fixing du 31 décembre couvre le 31 décembre, qui est en décembre. Les premiers jours de janvier
utilisent le fixing du premier jour ouvré de janvier.
\textbf{10.} À partir de contrats ou de dépôts à cheval sur la fin d'année et d'autres qui
ne le font pas, comme dans la \cref{prop:m2:money-markets:turn}~; ou à partir de l'historique des
fixings de fin d'année passés, utilisé comme a priori.
\textbf{11.} En cas de maintien, la moyenne de décembre vaut 3.8975~\%, le cours final 96.1025~:
$+6.25$ points de base, $6.25 \times 41.67 \times 1\,000 = +\text{USD}~260\,438$.
\textbf{12.} En cas de hausse, la moyenne vaut $(9 \times 3.8975 + 22 \times 4.1475)/31 =
4.0749\%$, cours 95.9251~: $-11.49$ points de base, $-\text{USD}~478\,869$.
\textbf{13.} $0.648 \times 260\,438 - 0.352 \times 478\,869 = 0$, au dollar près~: la
probabilité du marché est celle pour laquelle la transaction est équitable.
\textbf{14.} À 35.2~\%, parce que le cours est exactement la moyenne pondérée
par les probabilités des deux règlements-livraisons~; le desk ne gagne en moyenne que si sa
probabilité est inférieure à celle du marché.
\textbf{15.} Le spread du SOFR par rapport au taux directeur (conditions du repo, réserves),
passages de fin de mois et de fin d'année, et les fixings réalisés des jours qui précèdent la
réunion.
\textbf{16.} Les issues ne se limitent pas à 0 et à 25~: un mouvement de 50 points de base ou un
mouvement hors réunion est possible~; le cours intègre toutes les issues, si bien que la
répartition en deux issues n'est qu'une convention.
\textbf{17.} Les cours des contrats à terme intègrent des primes de risque~: les investisseurs paient pour se couvrir contre une hausse, si bien que
la probabilité implicite dans le cours peut dépasser la probabilité réelle~; et les
projections du Comité sont des prévisions, et non des paris.
\textbf{18.} Le taux applicable à partir de décembre serait d'environ $3.87 + 0.25 = 4.12\%$~; le cours de décembre devrait être relu à cette aune,
et les marchés réévaluent généralement l'ensemble
de la trajectoire après une surprise.
\textbf{19.} Résultat à retenir~: \emph{la probabilité implicite de hausse de décembre} est de 35.2~\% sans
\omterm{def:m2:money-markets:turn}{passage de fin de période} et de 33.4~\% avec un passage de 10 points de base au 31 décembre.
\textbf{20.} La moyenne du taux au jour le jour sur un mois, et donc chaque
décision et chaque \omterm{def:m2:money-markets:turn}{passage de fin de période} qui s'y situe.
\end{solution}

\begin{solution}{iq:m2:short-term-interest-rate-futures:1}
Prendre le contrat couvrant le mois de la réunion (ou mieux, le mois suivant si
la réunion est tardive)~; déduire le taux avant la réunion du contrat du mois
précédent~; résoudre la moyenne du mois pour trouver le taux postérieur à la réunion~; diviser
la variation par la taille du mouvement. Retirer d'abord les \omterm{def:m2:money-markets:turn}{passages de fin de période} et
le spread SOFR--taux directeur, et garder à l'esprit que c'est une approximation neutre au risque à deux issues.
\iqlookfor{la pondération par jour, le problème de la réunion tardive, et les mises en garde.}
\end{solution}

\begin{solution}{iq:m2:short-term-interest-rate-futures:2}
Afin que le cours évolue comme un cours obligataire~: il monte quand les taux baissent. Acheter le
contrat à terme revient alors à prendre une position longue sur des actifs sensibles aux taux, comme acheter un \omterm{def:m2:money-markets:bill}{bon du Trésor},
et les bourses peuvent afficher les cours, les cours acheteurs et les cours vendeurs dans le sens usuel~; une
valeur fixe par point de base fait du risque le nombre de contrats.
\iqlookfor{la convention de type obligataire et la valeur fixe par point de base.}
\end{solution}

\begin{solution}{iq:m2:short-term-interest-rate-futures:3}
Les contrats à terme sont réglés chaque jour, les forwards une seule fois à l'échéance. Un contrat à terme vendeur gagne lorsque les taux montent,
et réinvestit ces gains à des taux élevés~; il perd lorsque les taux baissent et finance ces pertes à des taux bas. Cette corrélation avantage le vendeur,
si bien que le taux du contrat à terme dépasse le taux forward~; dans un modèle gaussien, l'écart vaut environ
$\frac12\sigma^2 T_1 T_2$, négligeable à court terme,
et de quelques dizaines de points de base à dix ans.
\iqlookfor{le sens de l'écart et sa raison, ainsi que la croissance avec l'horizon.}
\end{solution}

\begin{solution}{iq:m2:short-term-interest-rate-futures:4}
Un pack regroupe quatre contrats trimestriels consécutifs négociés ensemble~: une année
du \omterm{def:m2:short-term-interest-rate-futures:strip}{strip de contrats à terme}. Le pack de la deuxième année isole les anticipations de taux à un an,
c'est-à-dire la trajectoire des taux directeurs après les prochaines réunions, sans l'exposition
des contrats les plus proches à la prochaine réunion et au bruit du marché monétaire~; c'est la position sur contrats à terme
la plus pure sur la trajectoire à moyen terme.
\iqlookfor{les packs comme moyen de négocier un segment de la courbe.}
\end{solution}

\begin{solution}{iq:m2:short-term-interest-rate-futures:5}
Fins de mois~: \omterm{def:m2:money-markets:turn}{passages de fin de période} dans les fixings, et week-ends qui font couvrir au dernier fixing
plusieurs jours. Réunions tardives~: le taux posté à la réunion n'est estimé qu'à partir de quelques jours, si bien que le bruit explose~;
il faut alors passer au contrat du mois suivant. Fin d'année~: un \omterm{def:m2:money-markets:turn}{passage de fin de période} important qui, s'il n'est pas retiré,
se lit comme une probabilité de mouvement. Il faut aussi tenir compte des cours obsolètes dans les mois peu liquides,
des calendriers de jours fériés
et du jour où le nouveau taux entre en vigueur.
\iqlookfor{des modes de défaillance concrets et leurs remèdes.}
\end{solution}

\begin{solution}{iq:m2:short-term-interest-rate-futures:6}
Non~: il s'agit d'une probabilité neutre au risque, qui contient la prime que les investisseurs paient pour
se couvrir contre l'issue qu'ils redoutent~; elle suppose aussi deux issues ainsi qu'un \omterm{def:m2:money-markets:turn}{passage de fin de période} et un spread connus.
Pour la convertir, il faut estimer la prime de risque (à partir de l'historique des décisions implicites
comparées aux décisions réalisées, ou à partir d'enquêtes), la comparer à la probabilité issue de son propre modèle,
et ne trader que l'écart net des coûts.
\iqlookfor{neutre au risque contre monde réel, et la manière d'estimer l'écart.}
\end{solution}
